\section{Problem Formulation}
\label{sec:problem_formulation}

We ask: \emph{Given a local wind disturbance, what is the minimum decision
scope that must be released to recover mission feasibility without
unnecessary replanning?} This short paper considers only wind-induced
travel-time and energy perturbations; no-fly zones, including
weather-induced ones, are excluded.
We use a \emph{hierarchical mixed
discrete--continuous optimization model}: assignment and visit order are
discrete, while paths and commodity transfers contain continuous decisions.
A dependency graph diagnoses the disturbance, and nested decision spaces
define the permitted releases. The minimum is measured within this declared
hierarchy. Detailed evaluator definitions and proofs are in the companion
technical supplement.

\subsection{Mission State and Remaining Decisions}
\label{subsec:information_remaining}

Let $\mathcal V=\{1,\ldots,N\}$ and $\mathcal U=\{1,\ldots,M\}$ denote
relief tasks and UAVs, respectively, with $N,M\ge1$.
The mission distributes one divisible commodity, with unit mass $m>0$.
The depot is $v_0$ at $\mathbf p_0$.
Task $i$ has location $\mathbf p_i$, demand $d_i\ge0$,
and mandatory quantity $d_i^{\mathrm{mand}}\in[0,d_i]$.
These quantities are fixed mission requirements; full-demand missions use
$d_i^{\mathrm{mand}}=d_i$.
All sites lie in a bounded polygonal operating region $\mathcal W$ of
positive diameter.

Fix an update epoch $t=t_{n+1}\in[0,T)$, where $T>0$ is the mission
horizon; epoch subscripts are suppressed below. Available information is
\begin{equation}
\mathscr I=(t,H,X^{\mathrm{fleet}},\widehat{\mathcal E}_{0:n+1},
\bar\Pi,\mathcal C^{\mathrm{fix}};\vartheta).
\label{eq:information_state}
\end{equation}
Here $H$ is executed history; $X^{\mathrm{fleet}}$ contains each UAV's
position, battery, onboard commodity quantity, and execution mode;
$\widehat{\mathcal E}_{0:n+1}$ is the received wind-packet history, with latest
packet $\widehat{\mathcal E}:=\widehat{\mathcal E}_{n+1}$;
$\bar\Pi$ is the incumbent;
and $\mathcal C^{\mathrm{fix}}$ specifies commitments. Fixed mission,
vehicle, operation, and representation data are collected in $\vartheta$.
History supplies consistent cumulative deliveries
$0\le q_i^{\mathrm{del}}(t)\le d_i$ and residual depot stock
$S^{\mathrm{dep}}\ge0$. Define
\begin{equation}
d_i^{\mathrm{rem}}=d_i-q_i^{\mathrm{del}}(t),\qquad
\underline d_i=[d_i^{\mathrm{mand}}-q_i^{\mathrm{del}}(t)]_+,
\label{eq:residual_requirements}
\end{equation}
where $[z]_+=\max\{z,0\}$, and let
$\mathcal V^{\mathrm{rem}}=\{i:d_i^{\mathrm{rem}}>0\}$.
Decisions use only received information and executed history.

A fixed finite spatial partition $(\mathcal W_r)_{r\in\mathcal R}$ and
time bins $(I_j)_{j\in\mathcal J}$ cover $\mathcal W\times[t,T]$ with unique
boundary ownership. Wind is constant on each cell--time bin.
The previous and received tables satisfy
\begin{equation}
\begin{split}
\widehat{\mathbf w}_{rj}^{+}
&=\widehat{\mathbf w}_{rj}^{-}+\delta\mathbf w_{rj},\\
\delta\mathbf w_{rj}&=\mathbf0\quad
\text{for }(r,j)\notin\mathcal B_\Delta,
\end{split}
\label{eq:local_wind_update}
\end{equation}
where $\mathcal B_\Delta\subsetneq\mathcal R\times\mathcal J$ is the
declared local support. Both tables obey
$\|\widehat{\mathbf w}_{rj}^{\omega}\|_2\le\bar w<\min_u v_u$,
$\omega\in\{-,+\}$, where $v_u>0$ is airspeed.
Wind changes travel time and energy; all spatial cells and directions
remain traversable.
Geometry, demands, stock, observed state, commitments, and operation rules
are identical in the paired evaluations; wind covers the entire remaining
horizon.

A remaining plan is $\Pi=(\mathcal A,\mathcal S,\mathcal P,\mathcal L)$.
Assignment $\mathcal A=(a_i)_{i\in\mathcal V^{\mathrm{rem}}}$ has
$a_i\in\mathcal U\cup\{\varnothing\}$: each visited task has one responsible
UAV, and $a_i=\varnothing$ exactly when no visit is planned. Repeated partial
deliveries by the responsible UAV are allowed.
Sequence $\mathcal S=(s_{u,h})$ uses fixed integer slot bounds $\bar n_u>0$,
with $s_{u,h}\in\mathcal V^{\mathrm{rem}}\cup\{v_0,\partial\}$ and
$\partial$ denoting inactive padding.
Its $n_u\in\{0,\ldots,\bar n_u\}$ active stops form a prefix following the
observed virtual origin $s_{u,0}$ and, when nonempty, end at the depot.
Intermediate depot visits
are allowed; a zero-slot route is permitted only for an idle UAV at the depot.
Paths $\mathcal P=(P_{u,h})$ are directed piecewise-linear curves in
$\mathcal W$ joining consecutive sites, with normalized arclength
parameterization. Quantities $\mathcal L=(x_{u,h},y_{u,h})$ specify
delivery and loading at stop $h$, including initial loading at $h=0$.
Inactive slots have canonical padding and zero transfers.

Let $\mathcal Y$ be the common domain enforcing these structural rules,
the observed starting state, and all commitments. For component
$a\in\{A,S,P,L\}$, fixed coordinate sets partition into committed $C^a$
and mutable $M^a$; with $\rho_B$ denoting projection onto coordinates $B$,
\begin{equation}
\rho_{C^a}(\Pi^a)=\rho_{C^a}(\bar\Pi^a),\qquad a\in\{A,S,P,L\}.
\label{eq:commitment_projection}
\end{equation}
Assume $\bar\Pi\in\mathcal Y$. Unfinished committed operations retain only
their residual durations and transfers; executed actions are not repeated.

\subsection{Execution and Mission Feasibility}
\label{subsec:hard_feasibility}

A given deterministic evaluator
$\Xi(\Pi)=\Gamma(\Pi;\mathscr I)$ induces times, payloads, batteries, and
flight-completion statuses from the observed state. Split each path at
wind-cell boundaries. On a resulting straight segment $e$ of length
$\ell_e$, entered at time $a$, let $w_{\parallel,e}$ and $w_{\perp,e}$
be the received wind components along and normal to its direction.
The constant-airspeed track model and linear flight-energy surrogate give
\begin{equation}
\begin{split}
s_{u,e}(s)&=w_{\parallel,e}(s)
+\sqrt{v_u^2-w_{\perp,e}(s)^2},\\
\int_a^{a+\tau_{u,e}}s_{u,e}(s)\,\mathrm ds&=\ell_e,\\
E_{u,e}&=\gamma_u(m_u^0+m_{u,e}^{\mathrm{pay}})v_u\tau_{u,e}.
\end{split}
\label{eq:wind_flight_model}
\end{equation}
Here $m_u^0>0$ is empty mass and $\gamma_u>0$ has units of energy per
mass per distance. On a segment of leg $h$, the onboard mass is
$m_{u,e}^{\mathrm{pay}}=m q_{u,h-1}^{+}$, where $q_{u,h-1}^{+}$ is the
commodity quantity after the preceding stop.
Ground speed is at least $v_u-\bar w>0$; a traversal is late if its
distance cannot be completed by $T$. Ground-operation rules are prescribed.
At depot departure, preparation loads supplies and restores the battery to
capacity $B_u>0$; elsewhere no energy is replenished. There is no independent
waiting decision or implicit depot queue. Ground operations occur at
assumed-safe sites.

Let $q_{u,h}^{-}$ and $q_{u,h}^{+}$ denote quantities before and after
stop $h$, initialized by the observed onboard load. Delivery occurs only
at relief visits, loading only at a depot before a subsequent leg, and
$x_{u,0}=y_{u,n_u}=0$. For active indices, feasibility requires
\begin{align}
q_{u,h}^{+}&=q_{u,h}^{-}-x_{u,h}+y_{u,h},
\; q_{u,h+1}^{-}=q_{u,h}^{+},
\label{eq:payload_flow}\\
0\le x_{u,h}&\le q_{u,h}^{-},\qquad
m q_{u,h}^{\pm}\le C_u^{\mathrm{pay}},
\label{eq:payload_capacity}\\
\sum_u\sum_{h=0}^{n_u-1}y_{u,h}&\le S^{\mathrm{dep}},
\label{eq:inventory}\\
Q_i(\Pi)&:=\sum_{u,h:\,s_{u,h}=i}x_{u,h},
\; \underline d_i\le Q_i(\Pi)\le d_i^{\mathrm{rem}},
\label{eq:service_bounds}
\end{align}
where $C_u^{\mathrm{pay}}>0$ is payload capacity and $y_{u,h}\ge0$.
The flight-flow equality applies only for $h<n_u$; empty sums are zero.
Returned payload remains onboard and is not restored to depot stock.
Battery states must remain in $[0,B_u]$, and each depot arrival after
flight must retain reserve $\beta_u\in[0,B_u]$ before recharge.
Positive deliveries to a task with deadline $\bar t_i$ must complete by
$\bar t_i$. Every UAV must finish at the depot by $T$ with at least
$\beta_u$. If a finite depot-bay capacity is specified, overlapping
preparation intervals must respect it.

Let $z=(\mathbf x,t_z,b_z,q_z)$ be a path's entering state.
The indicator $\operatorname{Comp}_u(P,z;\widehat{\mathcal E})$ equals one
exactly when $P$ completes by $T$; it does not certify battery sufficiency.

A reserve witness carries the current payload to the
depot without intermediate delivery or recharge. Its feasible set is
\begin{equation}
\begin{split}
\mathscr W_u(z)=\{W:\;&W(0)=\mathbf x,\ W(1)=\mathbf p_0,\\
&\operatorname{Comp}_u(W,z;\widehat{\mathcal E})=1,\\
&t_z+\tau_u(W,z)\le T,\quad
E_u(W,z)+\beta_u\le b_z\},
\end{split}
\label{eq:reserve_witness_set}
\end{equation}
where $W$ follows the same path convention and $\tau_u,E_u$ use the same
evaluator and packet. Require $\mathscr W_u(z)\ne\varnothing$ at each
initial state outside the depot, each initial depot departure, and every
post-service relief state. Zero-slot routes satisfy the terminal condition
directly. A failed supplied return path does not disprove existence.

Let $\mathcal Q^{\mathrm{hard}}$ collect the resource, service, deadline,
bay, reserve, and terminal predicates above. A complete evaluation has all
required nominal records defined and finishes all nominal operations by $T$.
A flight or ground operation that cannot
finish by $T$ violates the horizon; unavailable downstream values are $\bot$.
Predicates needing unavailable values are undetermined.
With $z_{u,h-1}^{\mathrm{dep}}$ the induced departure state, define
\begin{equation}
\begin{split}
\mathcal F(\mathscr I)=\{\Pi\in\mathcal Y:\;&
\operatorname{Complete}(\Xi(\Pi))=1,\\
&\operatorname{Comp}_u(P_{u,h},z_{u,h-1}^{\mathrm{dep}};
\widehat{\mathcal E})=1\ \ \forall(u,h),\\
&g(\Pi)=1\ \ \forall g\in\mathcal Q^{\mathrm{hard}}\}.
\end{split}
\label{eq:hard_feasible_set}
\end{equation}
Here $(u,h)$ ranges over active legs. Assume the incumbent is feasible
under the previous wind table at the same observed state,
$\bar\Pi\in\mathcal F^{-}$, where $\mathcal F^{-}$ replaces only the wind
in $\mathcal F(\mathscr I)$. Recovery is evaluated under the received table;
its certificate makes no assertion about future wind updates.

\subsection{Disturbance Support and Task Coupling}
\label{subsec:propagation_diagnostic}

Evaluate $\bar\Pi$ under the previous and received wind tables from the same
observed state. Unroll these paired evaluations into a finite typed
dependency DAG $\mathcal G$ containing all data and control dependencies,
including nominal flight, fixed diagnostic return paths, initial reserve
checks, wind queries, and time-bin selection guards.
Let $\mathcal N^\Delta$ contain the changed wind entries,
and let $\mathcal N^{\mathrm{seed}}$ contain the relevant atoms and their
direct wind-query children. Non-change roots agree between
projections; irrelevant atoms are isolated. The structural closure and
reportable support are
\begin{align}
\mathcal N^{\mathrm{cl}}&=\operatorname{Reach}_{\mathcal G}
(\mathcal N^{\mathrm{seed}}),
\label{eq:structural_closure}\\
\mathcal N^{\mathrm{rep}}&=\mathcal N^{\mathrm{seed}}\cup
\{v\in\mathcal N^{\mathrm{cl}}\setminus\mathcal N^\Delta:\mu(v)=1\}.
\label{eq:reportable_support}
\end{align}
Reachability includes its seeds. Materiality $\mu$ is one if exactly one
projected value is $\bot$, zero if both are $\bot$, and otherwise indicates
a difference exceeding a declared metric tolerance. Categorical tolerances
are zero; tolerances do not prune graph dependencies. Under complete
dependency representation, non-change values outside the closure agree.

Let $\mathcal V^{\mathrm{aff}}$ contain the remaining tasks incident to
legs whose diagnostic records intersect $\mathcal N^{\mathrm{rep}}$.
The reassignment graph is constructed separately over the common domain
$\mathcal Y$. Let $\mathcal Q^{\mathrm{tc}}$ contain the individual hard
constraints, including active flight and ground-operation completion.
Constraints retain their resource granularity: one stock constraint,
per-stop resource checks, and one bay-capacity constraint when imposed.
Inactive slot instances are satisfied by convention.

% Begin balancing within the first column of the final formulation page.
\balance
Write $\mathcal M=\mathop{\dot\bigcup}_{a}(\{a\}\times M^a)$ for the
tagged mutable coordinates. For each $g\in\mathcal Q^{\mathrm{tc}}$,
let $\eta_g(\Pi)$ contain its activation flag, plan-dependent arguments,
and definedness statuses, after substituting induced-state recursions.
A coordinate support $D_g\subseteq\mathcal M$ must satisfy
\begin{equation}
\begin{split}
\rho_{D_g}(\Pi)=\rho_{D_g}(\Pi')
\quad\Longrightarrow\quad\eta_g(\Pi)=\eta_g(\Pi'),\\
\Pi,\Pi'\in\mathcal Y.
\end{split}
\label{eq:constraint_coordinate_support}
\end{equation}
Thus dependencies include branch guards and upstream states, even when
$g$ is slack on the incumbent. Attribute each assignment coordinate to
its task and each UAV's mutable route or transfer coordinate to all tasks
that can be assigned to that UAV in $\mathcal Y$. Denote this conservative
attribution by $\Lambda(d)$. The task support and undirected incidence
graph are
\begin{align}
T(g)&=\bigcup_{d\in D_g}\Lambda(d),
\label{eq:constraint_task_support}\\
\mathcal B&=(\mathcal V^{\mathrm{rem}}\mathbin{\dot\cup}
\mathcal Q^{\mathrm{tc}},\mathcal E^{\mathrm{tc}}),\qquad
\mathcal E^{\mathrm{tc}}=\{\{i,g\}:i\in T(g)\}.
\label{eq:incidence_graph}
\end{align}
Here $g$ ranges over $\mathcal Q^{\mathrm{tc}}$.
Its task projection $\mathcal G^{\mathrm{task}}$ joins distinct $i,j$
exactly when $\{i,j\}\subseteq T(g)$ for some $g$.
The supplement specifies the attribution and a conservative construction
of $D_g$. These rules use only $\mathscr I$ and are fixed across levels.

Starting with $U_0=\mathcal V^{\mathrm{aff}}$, form the finite closure
\begin{equation}
U_{\ell+1}=U_\ell\cup
\bigcup_{g:\,T(g)\cap U_\ell\ne\varnothing}T(g),
\qquad U_\infty=\bigcup_{\ell\ge0}U_\ell.
\label{eq:task_closure_iteration}
\end{equation}
The Level-3 reassignment neighborhood is
\begin{equation}
\mathcal V^{R}=U_\infty\cap M^A,
\qquad U_\infty=\operatorname{Reach}_{\mathcal G^{\mathrm{task}}}
(\mathcal V^{\mathrm{aff}}).
\label{eq:reassignment_neighborhood}
\end{equation}
$U_\infty$ is the least superset of the affected tasks closed under every
declared constraint support: any support meeting it lies entirely within
it. Tasks with committed assignments remain in the graph and can connect
mutable tasks; the intersection with $M^A$ is taken only after closure.
This least-closure property does not assert a smallest feasible repair.
Shared inventory or unrestricted route participation can make
$\mathcal V^R=M^A$.

\subsection{Minimum Decision Release}
\label{subsec:minimum_escalation}

For $k\in\mathbb K:=\{0,1,2,3,4\}$, let $R_k^a\subseteq M^a$ be the
coordinates released at level $k$:
\begin{equation}
\begin{array}{@{}ccccc@{}}
\toprule
k&R_k^A&R_k^S&R_k^P&R_k^L\\
\midrule
0&\varnothing&\varnothing&\varnothing&\varnothing\\
1&\varnothing&\varnothing&M^P&\varnothing\\
2&\varnothing&M^S&M^P&M^L\\
3&\mathcal V^R&M^S&M^P&M^L\\
4&M^A&M^S&M^P&M^L\\
\bottomrule
\end{array}.
\label{eq:release_schedule}
\end{equation}
These levels permit retention, path repair, sequence and transfer repair,
neighborhood reassignment, and release of all \mbox{mutable} decisions.
The common commitment-preserving spaces are
\begin{equation}
\mathcal X_k=\{\Pi\in\mathcal Y:
\rho_{M^a\setminus R_k^a}(\Pi^a)
=\rho_{M^a\setminus R_k^a}(\bar\Pi^a)\ \ \forall a\}.
\label{eq:intervention_space}
\end{equation}
Cumulative release yields
$\{\bar\Pi\}=\mathcal X_0\subseteq\cdots\subseteq\mathcal X_4=\mathcal Y$.
Equality between adjacent spaces is allowed. Set
$\mathcal F_k=\mathcal X_k\cap\mathcal F(\mathscr I)$.
In particular, Level-3 has the explicit form
\begin{equation}
\mathcal X_3=\{\Pi\in\mathcal Y:
a_i=\bar a_i\quad\forall i\in M^A\setminus\mathcal V^R\}.
\label{eq:level_three_space}
\end{equation}
All mutable sequence, path, and transfer coordinates are already released
at Level-2; only assignment release is restricted by $\mathcal V^R$.
The full mission-feasibility test still applies to every candidate.

The minimum decision-release problem is
\begin{equation}
\begin{aligned}
k^F=\min_{k\in\mathbb K}\quad &k\\
\text{subject to}\quad &\exists\Pi\in\mathcal F_k,
\end{aligned}
\qquad \min\varnothing:=+\infty.
\label{eq:feasibility_escalation}
\end{equation}
Thus $k^F=+\infty$ exactly when $\mathcal F_4=\varnothing$.
For finite $k^F$, any witness
\begin{equation}
\Pi^F\in\mathcal F_{k^F}
\label{eq:feasible_selection}
\end{equation}
is an acceptable restored plan. Decision distance may optionally select
among feasible witnesses, as specified in the supplement; it imposes no
additional feasibility condition and cannot change $k^F$.

A validated witness in $\mathcal F_k$ proves $k^F\le k$.
Certifying equality also requires emptiness of all earlier levels;
failed search alone does not prove emptiness. Any certificate is relative
to the declared hierarchy, representation, commitments, and wind model.
Releasing a coordinate permits a change; it does not require one.

